%% 与 latex_render/templates/article.tex.j2 的导言区保持一致：
%% 构建期把这些宏包/类选项全部烤进 Tectonic 缓存，运行时无外网也能编译。
\documentclass[11pt,a4paper]{article}
\usepackage[utf8]{inputenc}
\usepackage[T1]{fontenc}
\usepackage{amsmath,amssymb}
\usepackage{graphicx}
\usepackage{booktabs}
\usepackage{longtable}
\usepackage{hyperref}
\usepackage{xcolor}
\usepackage[margin=2.5cm]{geometry}
\providecommand{\todo}[1]{\textcolor{red}{\textbf{[TODO: #1]}}}
\title{Warmup}\author{}\date{}
\begin{document}
\maketitle
\begin{abstract}warmup\end{abstract}
\section{Warmup}\label{sec:warmup}
Text with a citation~\cite{warmup2020sample} and \todo{placeholder}.
\begin{equation}E = mc^2\end{equation}
\begin{table}[t]\centering\begin{tabular}{ll}\toprule a & b \\ \midrule 1 & 2 \\ \bottomrule\end{tabular}\caption{t}\end{table}

%% longtable 预热：渲染器对**所有单栏模板**的表格都用 longtable（可分页、
%% 续页重复表头）。它一度不在任何预热文档里，于是 longtable.sty 从没进过
%% Tectonic 缓存——凡是带表格的论文，运行时都会 `File longtable.sty not found`，
%% 接着被降级逻辑注释掉宏包、再报 `Environment longtable undefined` 而彻底失败。
%% 预热文档的导言区必须与模板导言区逐条对齐，见 test_warmup_covers_templates。
\begin{longtable}{@{}p{0.5\linewidth}p{0.4\linewidth}@{}}
  \caption{longtable warmup}\label{tab:warm} \\
  \toprule
  {\raggedright\bfseries head a\par} & {\raggedright\bfseries head b\par} \\
  \midrule
  \endfirsthead
  \toprule
  {\raggedright\bfseries head a\par} & {\raggedright\bfseries head b\par} \\
  \midrule
  \endhead
  \midrule
  \endfoot
  \bottomrule
  \endlastfoot
  {\raggedright cell 1\par} & {\raggedright cell 2\par} \\
\end{longtable}

%% 字体族预热：数学下标/上标会拉取 cmr5/cmr7/cmmi7 等小号字，
%% 只有在构建期真正排版过一次，networkless 运行时才拿得到。
\section{Font warmup}
\begin{equation}\label{eq:warm}
  \sum_{i=1}^{n} x_i^{2} + \alpha_{jk} \approx \frac{1}{2}\int_0^\infty e^{-t^2}\,dt
\end{equation}
{\tiny tiny}\ {\scriptsize scriptsize}\ {\footnotesize footnotesize}\ {\small small}\
{\large large}\ {\Large Large}\ {\LARGE LARGE}\ {\huge huge}
\begin{itemize}\item item one\item item two\end{itemize}
\textbf{bold} \textit{italic} \texttt{mono} \underline{underline}\footnote{footnote warmup}


%% 字号 × 字形 全组合预热：Tectonic 按 *文件* 缓存字体，
%% \tiny 加粗这类组合会拉 lmroman6-bold.otf 等独立字体文件，
%% 只有构建期真的排版过，networkless + read-only 的运行时才拿得到。
{\tiny normal} {\tiny \textbf{bold}} {\tiny \textit{italic}} {\tiny \textsc{smallcaps}} {\tiny \texttt{mono}}

{\scriptsize normal} {\scriptsize \textbf{bold}} {\scriptsize \textit{italic}} {\scriptsize \textsc{smallcaps}} {\scriptsize \texttt{mono}}

{\footnotesize normal} {\footnotesize \textbf{bold}} {\footnotesize \textit{italic}} {\footnotesize \textsc{smallcaps}} {\footnotesize \texttt{mono}}

{\small normal} {\small \textbf{bold}} {\small \textit{italic}} {\small \textsc{smallcaps}} {\small \texttt{mono}}

{\normalsize normal} {\normalsize \textbf{bold}} {\normalsize \textit{italic}} {\normalsize \textsc{smallcaps}} {\normalsize \texttt{mono}}

{\large normal} {\large \textbf{bold}} {\large \textit{italic}} {\large \textsc{smallcaps}} {\large \texttt{mono}}

{\Large normal} {\Large \textbf{bold}} {\Large \textit{italic}} {\Large \textsc{smallcaps}} {\Large \texttt{mono}}

{\LARGE normal} {\LARGE \textbf{bold}} {\LARGE \textit{italic}} {\LARGE \textsc{smallcaps}} {\LARGE \texttt{mono}}

{\huge normal} {\huge \textbf{bold}} {\huge \textit{italic}} {\huge \textsc{smallcaps}} {\huge \texttt{mono}}

$x_{i_{j}}^{2^{3}} + \sqrt{\alpha_\beta}$

\bibliographystyle{plainnat}
\bibliography{refs}
\end{document}
